В данной главе рассматриваются особенности предметной области аналитических
бенчмарков, формулируются требования к современному генератору данных TPC-H и
выполняется сравнительный анализ существующих решений.

\subsection{TPC-H как стандарт аналитического бенчмаркинга}

\subsubsection{Назначение бенчмарка}

Бенчмарк TPC-H занимает устойчивое место в практике оценки аналитических СУБД,
поскольку задает одновременно структуру предметной области, правила
формирования набора данных и набор запросов, моделирующих типовые сценарии
аналитической обработки~\cite{tpch-spec}. В отличие от синтетических тестов,
оценивающих отдельные низкоуровневые операции, TPC-H ориентирован на проверку
сквозной эффективности систем поддержки принятия решений: выполнения
агрегирующих запросов, соединений, фильтрации и сортировки над большими
таблицами.

\subsubsection{Состав набора данных}

Схема TPC-H включает восемь взаимосвязанных таблиц~\cite{tpch-spec}. К ним
относятся region, nation, supplier,
part, partsupp, customer, orders и
lineitem. Для каждой таблицы спецификация задает состав атрибутов,
правила построения ключей, распределения значений и формулы вычисления
зависимых полей. Наличие строго определенных правил генерации делает TPC-H
удобной основой для воспроизводимых экспериментов: при одинаковом масштабе и
корректной реализации генератора результирующий набор данных должен быть
эквивалентен эталонному.

\subsubsection{Актуальность TPC-H в современных условиях}

Значение TPC-H усиливается тем, что он остается применимым и в современных
условиях, когда аналитические системы работают на многоядерных узлах,
используют параллельное выполнение и активно опираются на колонoчные методы
хранения и обработки данных~\cite{dewitt-gray}. Соответственно, требования к
генератору данных выходят за пределы простой записи строковых файлов и включают
поддержку масштабирования, эффективного внутреннего представления и
совместимости с современной экосистемой хранения данных.

\subsection{Требования современной аналитической инфраструктуры к генератору данных}

\subsubsection{Требования к выходному представлению данных}

Классический сценарий подготовки набора TPC-H предполагает формирование
строковых файлов, последующую конвертацию, разбиение на партиции и отдельную
загрузку в целевую систему~\cite{dbgen-mirror}. Такой конвейер увеличивает
число этапов обработки, требует дополнительного дискового пространства и
затрудняет проведение повторяемых экспериментов в случаях, когда необходимо
быстро получать наборы данных разного масштаба.

Современная аналитическая инфраструктура предъявляет к генератору данных более
широкий набор требований. Прежде всего данные должны формироваться в
представлении, согласованном с колонoчными моделями обработки, например в виде
батчей Apache Arrow~\cite{apache-arrow}. Кроме того, желательна прямая запись в
используемые на практике форматы хранения, прежде всего Parquet
~\cite{apache-parquet} и ORC~\cite{apache-orc}, а также в новые
форматы, ориентированные на интеграцию с аналитическими движками и
lakehouse-инфраструктурой, а также на работу с современными lakehouse-системами,
использующими табличные слои наподобие Apache Iceberg~\cite{apache-iceberg}.
К таким форматам относятся Lance~\cite{lance-format} и
Vortex~\cite{vortex-format}.

\subsubsection{Требования к размещению и доступу к данным}

Генератор должен учитывать особенности эксплуатации в распределенной
инфраструктуре. На практике результаты экспериментов нередко размещаются не
только на локальном диске, но и в объектных S3-совместимых
хранилищах. Это означает, что генератор должен уметь работать с абстракциями
файловых систем и URI, а не предполагать исключительно локальную
запись.

\subsubsection{Архитектурные требования}

Важным архитектурным требованием является отделение собственно генерации данных
от слоя сериализации. При таком разделении одна и та же логика формирования
таблиц может использоваться как для автономной выгрузки на носитель, так и для
встроенных сценариев, в которых данные передаются в оперативной памяти
непосредственно потребляющей системе. Этот подход особенно важен для
OLAP-движков, где чтение с диска в некоторых экспериментах целесообразно
заменять прямой генерацией колонoчных батчей.

\subsection{Анализ существующих решений}

\subsubsection{Генератор dbgen}

\paragraph{Достоинства}

Исторически основным средством подготовки набора TPC-H является генератор
dbgen, поставляемый вместе со спецификацией бенчмарка
TPC-H~\cite{tpch-spec,dbgen-mirror}. Его сильной стороной является полное
следование эталонным правилам генерации и фактический статус исходной точки для
проверки корректности других реализаций.

\paragraph{Ограничения}

Вместе с тем dbgen отражает архитектурные предпосылки своего времени.
Он ориентирован на формирование строковых файлов, не использует современные
внутрипамятные колонoчные представления, не обеспечивает прямой записи в
современные форматы аналитического хранения и плохо приспособлен к сценариям,
где требуется интеграция с удаленными объектными хранилищами. Дополнительно
практическое использование dbgen ограничивается условиями лицензии TPC,
что затрудняет его прямую адаптацию для некоторых исследовательских и
коммерческих сценариев.

\subsubsection{Генератор tpchgen-rs}

\paragraph{Достоинства}

Одним из современных открытых аналогов является tpchgen-rs
~\cite{tpchgen-rs}. Это решение реализовано на языке Rust и ориентировано на
параллельную генерацию таблиц и партиций. Существенным достоинством
tpchgen-rs является поддержка прямого формирования данных в
представлении Apache Arrow~\cite{apache-arrow}, что делает его существенно ближе к современной
аналитической инфраструктуре, чем dbgen.

\paragraph{Ограничения}

При этом tpchgen-rs в основном ориентирован на локальные сценарии
использования и не решает в полной мере задачу прямой записи в удаленные
S3-совместимые хранилища. Кроме того, для него не является центральной
задача архитектурного отделения ядра генерации от слоя сериализации таким
образом, чтобы генератор можно было использовать как универсальный встраиваемый
источник данных для аналитических систем.

\subsubsection{Сравнительный анализ и обоснование направления разработки}

Сравнение рассмотренных решений показывает, что ни одно из них не сочетает
одновременно все свойства, значимые для современного сценария эксплуатации
аналитических систем: прямую генерацию в колонoчное представление, запись в
современные форматы хранения, работу с S3-совместимыми хранилищами,
параллельное выполнение и возможность непосредственного встраивания в
аналитические движки.

Именно эта совокупность требований определяет направление разработки
Schengen~\cite{schengen-repo}. Генератор должен сохранять
детерминированность TPC-H и при этом быть ориентированным не на промежуточные
строковые файлы, а на прямое формирование батчей Apache Arrow~\cite{apache-arrow} с последующей
сериализацией в актуальные форматы. Дополнительное расширение области
применения дает интеграция генератора с DuckDB~\cite{schengen-duckdb-repo} в виде отдельного расширения,
использующего прямую подачу данных в память СУБД.

\begin{table}[h]
\caption{Сравнение генераторов данных TPC-H}
\label{tab:analogs}
\centering
\small
\begin{tabular}{|p{0.28\textwidth}|p{0.19\textwidth}|p{0.19\textwidth}|p{0.19\textwidth}|}
\hline
\textbf{Критерий} & \textbf{dbgen} & \textbf{tpchgen-rs} & \textbf{Schengen} \\
\hline
Основной формат вывода & tbl & прямая генерация, Arrow & прямая генерация, Arrow \\
\hline
Современные форматы хранения & требуется отдельная конвертация & частично поддерживаются & поддерживаются напрямую \\
\hline
Параллельная генерация & отсутствует & поддерживается & поддерживается \\
\hline
Работа с S3-совместимыми хранилищами & отсутствует & не является базовым сценарием & поддерживается напрямую \\
\hline
Встраивание в аналитические системы & практически отсутствует & ограничено & предусмотрено архитектурой \\
\hline
Лицензионная модель & ограничения TPC & открытая & открытая \\
\hline
\end{tabular}
\end{table}

\subsection*{Выводы по главе}
\addcontentsline{toc}{subsection}{Выводы по главе}

В главе показано, что TPC-H остается актуальной основой для оценки
аналитических СУБД, однако требования к генераторам данных в современной
инфраструктуре существенно расширились. Помимо корректного воспроизведения
эталонной схемы и распределений, генератор должен поддерживать колонoчное
внутрипамятное представление, прямую сериализацию в современные форматы,
работу с удаленными хранилищами и встраивание в аналитические системы.

Анализ dbgen и tpchgen-rs показывает, что существующие
решения закрывают эти требования лишь частично. Это обосновывает разработку
Schengen как генератора, ориентированного на современные сценарии
аналитической эксплуатации при сохранении детерминированности и совместимости с
правилами TPC-H.
